\section{Approach}\label{sec:Approach}

This section presents our approach to detecting missing \texttt{RB\_GC\_GUARD} directives. We first discuss our choice of analysis framework and the rationale for selecting CodeQL (Section~\ref{sec:tool_selection}). We then describe the detection algorithm and its key predicates (Section~\ref{sec:detection_algorithm}). Finally, we detail the implementation of each analysis component (Section~\ref{sec:implementation}).

\subsection{Tool Selection}\label{sec:tool_selection}

Given the challenges outlined in Section~\ref{sec:Challenges}, we considered several approaches for implementing the static analysis. The ideal tool must support efficient querying over a large codebase, provide access to both syntactic (AST) and semantic (control-flow, type) information, enable recursive queries for computing transitive closures, and allow declarative specification of complex patterns.

Traditional approaches to static analysis include custom analyzers written in general-purpose languages, pattern-matching tools such as Coccinelle\cite{UgawaFujimoto2019PRO}, and compiler-integrated analyses. Custom analyzers offer maximum flexibility but require substantial engineering effort to parse C code, construct control-flow graphs, and implement efficient indexing. Pattern-matching tools excel at syntactic transformations but lack the semantic analysis needed to track control-flow ordering and compute call-graph reachability. Compiler-integrated analyses (e.g., GCC plugins, Clang checkers) provide rich semantic information but are tightly coupled to specific compiler versions and optimization pipelines.

We selected CodeQL as our implementation platform for several reasons. First, CodeQL represents source code as a queryable relational database, with pre-computed indices over types, names, and structural properties\cite{ql_ecoop}. This enables efficient lookups that avoid the need for full enumeration. Second, CodeQL provides built-in support for C/C++ analysis, including AST matching, control-flow analysis functionability, and type inference. Third, the QL language supports recursive predicates \cite{codeql_recursion}, which is essential for computing transitive call-graph reachability. Fourth, CodeQL's declarative query language allows patterns to be specified at a high level of abstraction while the query optimizer handles efficient execution.

These capabilities align well with the challenges identified in Section~\ref{sec:Challenges}: indexed database access addresses scale (C1), built-in macro and type support addresses derivation diversity (C2), recursive queries address GC trigger  (C3), macro metadata addresses guard detection (C4).

\subsection{Query Implementation}\label{sec:detection_algorithm}

Our detection algorithm implements the query requirements from Section~\ref{sec:Challenges} as a conjunction of predicates over the CodeQL database. The core detection pattern can be expressed as an ordering constraint over control-flow paths:

\[
\textsc{Derivation}(p \leftarrow v) \xrightarrow{CFG^*} \textsc{LastUse}(v) \xrightarrow{CFG^+} \textsc{Trigger} \xrightarrow{CFG^+} \textsc{Use}(p)
\]

where $\xrightarrow{CFG^*}$ denotes zero or more control-flow steps and $\xrightarrow{CFG^+}$ denotes one or more steps. The constraint captures that the last \texttt{VALUE}-level use of $v$ must precede the trigger, which in turn must precede the use of $p$.

Listing~\ref{lst:running_example} illustrates a buggy pattern. The local variable \texttt{str} is used to derive a pointer \texttt{ptr} on line 3. The last \texttt{VALUE}-level access to \texttt{str} occurs on line 4 via \texttt{RSTRING\_LEN}. A GC-triggering call occurs on line 6 (\texttt{rb\_str\_new}), after which \texttt{ptr} is used on line 8. Because \texttt{str} has no access after the trigger and no explicit guard, an optimizing compiler may treat \texttt{str} as dead at line 6, potentially causing \texttt{ptr} to become a dangling pointer.

\begin{figure}[tb]
\begin{lstlisting}[caption={Running example: a buggy pattern in CRuby extension code.},label={lst:running_example}]
static VALUE process_data(VALUE input) {
    VALUE str = rb_funcall(input, rb_intern("to_s"), 0);
    char *ptr = RSTRING_PTR(str);     // derivation
    size_t len = RSTRING_LEN(str);    // last VALUE-level use

    VALUE result = rb_str_new(0, len * 2);  // GC trigger

    memcpy(RSTRING_PTR(result), ptr, len);  // use of ptr
    return result;
}
\end{lstlisting}
\end{figure}

The query proceeds in four phases. In the first phase, we identify all local variables of type \texttt{VALUE}. In the second phase, for each \texttt{VALUE} variable, we find derivation sites where a raw pointer is derived. In the third phase, we compute the set of functions that may transitively trigger GC via call-graph reachability. In the fourth phase, we check control-flow ordering constraints to identify variables that need guards but lack them.

\subsection{Implementation}\label{sec:implementation}

We implemented GuardLint as a set of CodeQL queries distributed as a query pack.\footnote{Artifact: \url{https://github.com/dummyx/guardql/tree/e8e4ed62f5ec89bd663ac380dccd2c933a48186d}. The main queries used in this paper are \texttt{missing\_guards\_vo.ql}, \texttt{redundant\_guards.ql}, and \texttt{good\_guards.ql} under \texttt{codeql-custom-queries-cpp}.}

\subsubsection*{Identifying VALUE Variables}
The foundation of our analysis is identifying local variables of type \texttt{VALUE}. Listing~\ref{lst:value_variable_detection} shows the CodeQL class that captures this. The predicate handles both direct type declarations and typedef indirections.

\begin{figure}[tb]
\begin{lstlisting}[caption={CodeQL class for identifying \texttt{VALUE} variables.},label={lst:value_variable_detection}]
class ValueVariable extends LocalVariable {
  ValueVariable() {
    this.getType().getName() = "VALUE"
    or
    this.getType().(TypedefType).getName() = "VALUE"
  }
}
\end{lstlisting}
\end{figure}

\subsubsection*{Detecting Pointer Derivations}
We define predicates that recognize common derivation operations in CRuby. Listing~\ref{lst:derivation_predicate} shows the structure of our derivation detection. We try to detect this by matching the underlying function call's function name.

\begin{figure}[tb]
\begin{lstlisting}[caption={CodeQL predicate for detecting pointer derivations (simplified).},label={lst:derivation_predicate}]
class InnerPointerTakingFunctionByName extends Function {
  InnerPointerTakingFunctionByName() {
    this.getName() in [
        "RSTRING_PTR"
        // Dedacted
      ]
  }
}
\end{lstlisting}
\end{figure}

This is necessarily an over-approximation: not all recognized APIs always produce unsafe pointers, and some unsafe derivations may be missed if they use unmodeled patterns. However, the modeled set covers the most common and well-documented APIs in CRuby's extension interface.

\subsubsection*{Computing GC Trigger Reachability}
We identify GC-triggering functions through transitive reachability to CRuby's GC entry point \texttt{gc\_enter}. This is achieved using CodeQL's recursive predicate capability\cite{codeql_recursion}. Listing~\ref{lst:is_gc_trigger} shows the implementation.

\begin{figure}[tb]
\begin{lstlisting}[caption={CodeQL predicate for GC trigger detection via call-graph reachability.},label={lst:is_gc_trigger}]
predicate isGcTrigger(Function function) {
  exists(Expr s, Call call |
    s = function.getAPredecessor*() and
    s instanceof FunctionCall and
    s.getAChild*() = call and
    (call.getTarget().getName() = "gc_enter" or
     isGcTrigger(call.getTarget()))
  )
}

class GcTriggerCall extends FunctionCall {
  GcTriggerCall() {
    isGcTrigger(this.getTarget())
  }
}
\end{lstlisting}
\end{figure}

This predicate computes the transitive closure of the call graph rooted at \texttt{gc\_enter}. It is a may-trigger approximation that intentionally errs on the side of reporting candidates, addressing the challenge that most CRuby functions can transitively reach GC. The recursive formulation leverages CodeQL's built-in memoization and fixed-point computation, making it efficient even for large call graphs. The result is cached in the database, so subsequent queries can reuse the computed trigger set.

\subsubsection*{Detecting Existing Guards}
GuardLint detects \texttt{RB\_GC\_GUARD} usage through recognizing the characteristic expansion pattern of the GCC-style implementation: the introduction of a local temporary named \texttt{rb\_gc\_guarded\_ptr} initialized with the address of the guarded \texttt{VALUE}.

\begin{figure}[tb]
\begin{lstlisting}[caption={CodeQL predicate for detecting existing guards.},label={lst:guard_detection}]
predicate hasGuard(ValueVariable v) {
  // Primary: macro invocation metadata
  exists(MacroInvocation mi |
    mi.getMacroName() = "RB_GC_GUARD" and
    mi.getExpr().getAChild*().(VariableAccess)
      .getTarget() = v
  )
  or
  // Fallback: characteristic expansion pattern
  exists(LocalVariable guardPtr, AddressOfExpr addr |
    guardPtr.getName() = "rb_gc_guarded_ptr" and
    guardPtr.getInitializer().getExpr() = addr and
    addr.getOperand().(VariableAccess).getTarget() = v
  )
}
\end{lstlisting}
\end{figure}

\subsubsection*{Control-Flow Ordering Analysis}
The core detection logic combines the above components with control-flow ordering constraints. We compute the last syntactic \texttt{VALUE}-level access to the owning variable along control-flow paths, using CodeQL's control-flow relations\cite{codeql_cpp_getapredecessor}. Listing~\ref{lst:ordering_constraints} shows the key predicates for establishing temporal ordering.

\begin{figure}[tb]
\begin{lstlisting}[caption={CodeQL predicates for control-flow ordering constraints.},label={lst:ordering_constraints}]
// Check if expr1 precedes expr2 in CFG
predicate precedes(Expr expr1, Expr expr2) {
  expr1.getASuccessor+() = expr2
}

// Find last VALUE-level access to v before a given point
predicate isLastValueAccess(VariableAccess va,
                            ValueVariable v,
                            Expr beforePoint) {
  va.getTarget() = v and
  precedes(va, beforePoint) and
  not exists(VariableAccess later |
    later.getTarget() = v and
    precedes(va, later) and
    precedes(later, beforePoint)
  )
}

// Core detection predicate
predicate needsGuard(ValueVariable v,
                     Variable p,
                     GcTriggerCall trigger,
                     Expr pointerUse,
                     Expr derivation) {
  isDerivation(v, derivation) and
  derivation.getASuccessor*() = trigger and
  trigger.getASuccessor+() = pointerUse and
  pointerUse.getAChild*().(VariableAccess).getTarget() = p and
  isLastValueAccess(_, v, trigger)
}
\end{lstlisting}
\end{figure}

A \texttt{VALUE} variable is flagged as needing a guard when a subordinate pointer $p$ is used after a modeled trigger site, the owning \texttt{VALUE} $v$ is not accessed after the trigger (so $v$ may be treated as dead by the compiler), and no explicit \texttt{RB\_GC\_GUARD} protects $v$ across the trigger.

\subsubsection*{Main Query Structure}
The main query structure combines these components. Listing~\ref{lst:missing_guard_ql} shows the missing guard detection query.

\begin{figure}[tb]
\begin{lstlisting}[caption={Main query for detecting missing guards.},label={lst:missing_guard_ql}]
/**
 * @name Missing RB_GC_GUARD
 * @description Finds VALUE variables that may need
 *              RB_GC_GUARD but lack one.
 * @kind problem
 * @problem.severity warning
 */
from ValueVariable v, Variable p,
     GcTriggerCall gtc,
     Expr pointerUsageAccess,
     Expr innerPointerTaking
where
  needsGuard(v, p, gtc, pointerUsageAccess,
             innerPointerTaking)
  and not hasGuard(v)
select v,
  "VALUE $@ may need RB_GC_GUARD: pointer derived at $@, " +
  "GC may trigger at $@, pointer used at $@.",
  v, v.getName(),
  innerPointerTaking, "here",
  gtc, "this call",
  pointerUsageAccess, "here"
\end{lstlisting}
\end{figure}

Because the CFG represents one possible evaluation order out of those allowed by the C language specification and compiler, the analysis remains a conservative approximation rather than a compiler-precise liveness proof. Additionally, we conservatively treat passing a subordinate pointer as an argument as a potential use, though we do not fully analyze derived-pointer dereferences inside callees in the current configuration.

\subsubsection*{Complementary Queries}
In addition to the missing-guard query, GuardLint includes queries to identify \emph{matched guards} (existing \texttt{RB\_GC\_GUARD} invocations where the guarded variable satisfies \texttt{needsGuard}, representing guards that our model confirms as necessary) and \emph{potentially redundant guards} (existing guards where the guarded variable does \emph{not} satisfy \texttt{needsGuard}, which may be candidates for removal, though they could also represent defensive coding or protection against future code changes).

\begin{figure}[tb]
\begin{lstlisting}[caption={Query for detecting potentially redundant guards.},label={lst:redundant_guard_ql}]
/**
 * @name Potentially redundant RB_GC_GUARD
 * @description Finds guards that may not be necessary
 *              under our model.
 * @kind problem
 * @problem.severity recommendation
 */
from ValueVariable v, MacroInvocation guard
where
  guard.getMacroName() = "RB_GC_GUARD" and
  guard.getExpr().getAChild*().(VariableAccess)
    .getTarget() = v and
  not needsGuard(v, _, _, _, _)
select guard,
  "RB_GC_GUARD for $@ may be redundant under our model.",
  v, v.getName()
\end{lstlisting}
\end{figure}

These complementary queries support a practical workflow: developers can review matched guards for confidence, investigate missing guards for safety, and consider removing redundant guards for performance.
